%%%%%%%%%%%%%%%%%%%%%%%%%%%%%%%%%%%%%%%%%%%%%%%%%%%%%%%%%%
% FILE: chapter09_guia_uso.tex
%
% PURPOSE
% -------------------------------------------------------
% Provides a step-by-step guide for using the template.
%
%
% ARCHITECTURE ROLE
% -------------------------------------------------------
% Practical user guide.
%
% This chapter explains how to start a project using
% the template and how to work with its main features.
%
%
% DEPENDENCIES
% -------------------------------------------------------
% Entire template structure
%
%
% NOTES FOR DEVELOPERS
% -------------------------------------------------------
% This chapter should remain procedural and explicit.
%
% The goal is to allow both humans and automated agents
% to understand how to operate the template correctly.
%
%%%%%%%%%%%%%%%%%%%%%%%%%%%%%%%%%%%%%%%%%%%%%%%%%%%%%%%%%%


\chapter{Guía de uso de la plantilla}

Este capítulo explica paso a paso cómo utilizar la
plantilla para escribir un documento académico.



% --------------------------------------------------
% PASO 1
% --------------------------------------------------

\section{Paso 1 — Obtener la plantilla}

El primer paso consiste en descargar la plantilla.

Esto puede hacerse de dos formas.

\begin{itemize}

\item descargando el archivo \texttt{.zip} del repositorio
\item clonando el repositorio con Git

\end{itemize}



\begin{lstlisting}[caption={Clonar repositorio},label={code:clone}]
git clone https://github.com/usuario/latex-template.git
\end{lstlisting}

\source{Plantilla}



Una vez descargada la plantilla, el usuario debe abrir
el archivo principal:

\texttt{main.tex}



% --------------------------------------------------
% PASO 2
% --------------------------------------------------

\section{Paso 2 — Configurar los metadatos}

Antes de escribir el contenido del documento es necesario
configurar los metadatos.

Estos se encuentran en el archivo:

\texttt{00\_config/config.tex}



Ejemplo:

\begin{lstlisting}[caption={Configuración de metadatos},label={code:metadata}]
\newcommand{\documenttitle}{Titulo del Trabajo}
\newcommand{\documentauthor}{Nombre del Autor}
\newcommand{\degree}{Grado en Ingenieria Informatica}
\newcommand{\faculty}{Facultad de Ciencias y Tecnologia}
\newcommand{\university}{UNIE Universidad}
\end{lstlisting}

\source{Plantilla}



\begin{notebox}
Estos metadatos se utilizan automáticamente en
varias partes del documento, como la portada
y los metadatos del PDF.
\end{notebox}



% --------------------------------------------------
% PASO 3
% --------------------------------------------------

\section{Paso 3 — Crear capítulos}

El contenido principal del documento se escribe en la
carpeta:

\texttt{02\_chapters}



Cada capítulo debe almacenarse en un archivo independiente.

Ejemplo:

\begin{lstlisting}[caption={Crear un capítulo},label={code:createchapter}]
\chapter{Introduccion}

Este capitulo introduce el tema del documento.
\end{lstlisting}

\source{Plantilla}



Después, el capítulo debe añadirse en el archivo
\texttt{main.tex}.

\begin{lstlisting}[caption={Cargar capítulo},label={code:loadchapter}]
%%%%%%%%%%%%%%%%%%%%%%%%%%%%%%%%%%%%%%%%%%%%%%%%%%%%%%%%%%
% FILE: chapter01_introduccion.tex
%
% PURPOSE
% -------------------------------------------------------
% Introduces the LaTeX template and explains its
% objectives, philosophy and intended use.
%
%
% ARCHITECTURE ROLE
% -------------------------------------------------------
% First chapter of the template documentation.
%
% Provides conceptual context before explaining
% the technical structure of the project.
%
%
% DEPENDENCIES
% -------------------------------------------------------
% environments.tex
% (definitionbox, notebox, warningbox)
%
%
% NOTES FOR DEVELOPERS
% -------------------------------------------------------
% This chapter should remain conceptual.
%
% Do not include implementation details here.
% Technical explanations belong to later chapters.
%
%%%%%%%%%%%%%%%%%%%%%%%%%%%%%%%%%%%%%%%%%%%%%%%%%%%%%%%%%%


\chapter{Introducción}

Esta plantilla ha sido diseñada para facilitar la redacción de
documentos académicos y técnicos utilizando \LaTeX.

El objetivo principal es proporcionar una estructura modular
que permita separar claramente:

\begin{itemize}
\item configuración del documento
\item estilos tipográficos
\item contenido académico
\item recursos gráficos y bibliográficos
\end{itemize}

Gracias a esta separación, el autor puede concentrarse
exclusivamente en escribir el contenido del documento
sin preocuparse por el formato.



% --------------------------------------------------
% OBJETIVO DE LA PLANTILLA
% --------------------------------------------------

\section{Objetivo de la plantilla}

\begin{definitionbox}[Plantilla académica]
Una plantilla académica es una estructura predefinida
que automatiza la configuración tipográfica y
organizativa de un documento.
\end{definitionbox}

Esta plantilla está diseñada para:

\begin{itemize}
\item trabajos universitarios
\item informes técnicos
\item proyectos de ingeniería
\item Trabajos Fin de Grado (TFG)
\item Trabajos Fin de Máster (TFM)
\end{itemize}



% --------------------------------------------------
% FILOSOFÍA DE DISEÑO
% --------------------------------------------------

\section{Filosofía de diseño}

La plantilla sigue varios principios de diseño inspirados
en prácticas de ingeniería de software.

\begin{itemize}
\item modularidad
\item separación de responsabilidades
\item reutilización
\item documentación interna
\end{itemize}



\begin{notebox}
El objetivo no es únicamente generar un documento
correctamente formateado, sino también demostrar
buenas prácticas en la organización de proyectos
\LaTeX.
\end{notebox}



% --------------------------------------------------
% ORGANIZACIÓN DEL PROYECTO
% --------------------------------------------------

\section{Organización del proyecto}

El proyecto se organiza en varias carpetas que separan
la configuración del documento del contenido académico.



\begin{lstlisting}[caption={Estructura general del proyecto},label={code:estructura}]
main.tex

latex-university-template/
|-- main.tex
|-- 00_config/
|   |-- config.tex
|   `-- packages.tex
|-- 01_frontmatter/
|-- 02_chapters/
|-- 03_figures/
|-- 04_bibliography/
|-- 05_appendices/
|-- 06_styles/
`-- 07_acronyms/
\end{lstlisting}

\source{Plantilla}



% --------------------------------------------------
% USO RECOMENDADO
% --------------------------------------------------

\section{Uso recomendado}

El flujo de trabajo recomendado es el siguiente:

\begin{enumerate}
\item Configurar los metadatos del documento
\item Escribir el contenido en la carpeta \texttt{02\_chapters}
\item Añadir figuras en \texttt{03\_figures}
\item Añadir referencias en \texttt{05\_bibliography}
\end{enumerate}



\begin{warningbox}
Modificar directamente los archivos de configuración
o de estilos sin comprender su función puede producir
errores de compilación o inconsistencias tipográficas.
\end{warningbox}



% --------------------------------------------------
% METODOLOGÍA DE DOCUMENTACIÓN
% --------------------------------------------------

\section{Metodología de documentación}

Todos los archivos de la plantilla siguen una convención
de comentarios estructurados.

Cada archivo comienza con un bloque de documentación
que describe su función dentro del proyecto.

\begin{lstlisting}[caption={Bloque de documentación estándar},label={code:docblock}]
%%%%%%%%%%%%%%%%%%%%%%%%%%%%%%%%%%%%%%%%%%%%%%%%%%%%%%%%%%
% FILE: filename.tex
%
% PURPOSE
% Describe el proposito del archivo
%
% ARCHITECTURE ROLE
% Describe su funcion dentro del sistema
%
% DEPENDENCIES
% Archivos o paquetes necesarios
%
% NOTES FOR DEVELOPERS
% Indicaciones para mantenimiento
%%%%%%%%%%%%%%%%%%%%%%%%%%%%%%%%%%%%%%%%%%%%%%%%%%%%%%%%%%
\end{lstlisting}

\source{Plantilla}



\begin{examplebox}
Este sistema de documentación permite comprender
rápidamente la arquitectura del proyecto al revisar
los archivos de la plantilla.
\end{examplebox}

\end{lstlisting}

\source{Plantilla}



% --------------------------------------------------
% PASO 4
% --------------------------------------------------

\section{Paso 4 — Añadir figuras}

Las imágenes deben almacenarse en la carpeta:

\texttt{03\_figures}



Las figuras se insertan mediante la macro:

\begin{lstlisting}[caption={Insertar figura},label={code:insertfigure}]
\figura{03_figures/imagen}{
Topologia de red
}{Elaboracion propia}{topologia}
\end{lstlisting}

\source{Plantilla}



% --------------------------------------------------
% PASO 5
% --------------------------------------------------

\section{Paso 5 — Añadir tablas}

Las tablas se insertan mediante la macro:

\begin{lstlisting}[caption={Insertar tabla},label={code:table}]
\tabla{
\begin{tabular}{cc}
Dato 1 & Dato 2 \\
Dato 3 & Dato 4
\end{tabular}
}{Ejemplo de tabla}{Plantilla}{tablaejemplo}
\end{lstlisting}

\source{Plantilla}



% --------------------------------------------------
% PASO 6
% --------------------------------------------------

\section{Paso 6 — Añadir código}

El código fuente se muestra mediante el entorno
\texttt{codeblock}.

\begin{lstlisting}[caption={Bloque de código},label={code:codeblock}]
\begin{codeblock}{Ejemplo de codigo}{ejemplo}
for i in range(10):
    print(i)
\end{codeblock}

\source{Plantilla}
\end{lstlisting}

\source{Plantilla}



% --------------------------------------------------
% PASO 7
% --------------------------------------------------

\section{Paso 7 — Añadir acrónimos}

Los acrónimos se definen en:

\texttt{07\_acronyms/acronyms.tex}

Ejemplo:

\begin{lstlisting}[caption={Definición de acrónimo},label={code:acronym}]
\DeclareAcronym{api}{
  short = API ,
  long  = Application Programming Interface
}
\end{lstlisting}

\source{Plantilla}



Uso en el documento:

\begin{lstlisting}[caption={Uso de acrónimo},label={code:acronymuse}]
\ac{api}
\end{lstlisting}

\source{Plantilla}



Ejemplo en texto:

\begin{examplebox}
Las \ac{api} permiten la comunicación entre
diferentes sistemas de software.
\end{examplebox}



% --------------------------------------------------
% PASO 8
% --------------------------------------------------

\section{Paso 8 — Añadir referencias}

Las referencias bibliográficas se almacenan en:

\texttt{04\_bibliography/references.bib}



Ejemplo de referencia:

\begin{lstlisting}[caption={Referencia BibTeX},label={code:bib}]
@book{knuth1984,
  author = {Donald E. Knuth},
  title = {The TeXbook},
  year = {1984},
  publisher = {Addison-Wesley}
}
\end{lstlisting}

\source{Plantilla}



Uso en el texto:

\begin{lstlisting}[caption={Citar referencia},label={code:cite}]
\cite{knuth1984}
\end{lstlisting}

\source{Plantilla}



Ejemplo:

\begin{examplebox}
El sistema tipográfico \TeX{} fue desarrollado por
Donald Knuth \cite{knuth1984}.
\end{examplebox}



% --------------------------------------------------
% PASO 9
% --------------------------------------------------

\section{Paso 9 — Compilar el documento}

Para generar el documento final es necesario compilar
el proyecto.

El proceso típico es el siguiente:

\begin{lstlisting}[caption={Proceso de compilación},label={code:compile}]
pdflatex main.tex
biber main
pdflatex main.tex
pdflatex main.tex
\end{lstlisting}

\source{Plantilla}



\begin{warningbox}
Si no se ejecuta \texttt{biber}, las referencias
bibliográficas no se generarán correctamente.
\end{warningbox}



% --------------------------------------------------
% RESULTADO
% --------------------------------------------------

\section{Resultado}

Después de completar estos pasos, el documento
final se generará como un archivo PDF completamente
formateado.



\begin{examplebox}
La estructura modular de la plantilla permite
concentrarse en el contenido del documento sin
preocuparse por los detalles tipográficos.
\end{examplebox}
